\section{Chương~\ref{sec:gaba}. \nt{GLUT} và \nt{GABA} cùng nhau}

Các khẳng định logic và động lực học của chương được suy ra ngay trong chương từ phân tích tuyến tính và định luật Ohm; thực nghiệm cho thấy các mạch làm nền cho chúng --- phép cấm có độ trễ, ức chế thuận theo sau kích thích, sự ổn định nhờ ức chế, nhịp của vòng \nt{GLUT}--\nt{GABA} --- thực sự hoạt động trong não.

{\footnotesize
\setlength{\tabcolsep}{3pt}\renewcommand{\arraystretch}{1.15}
\begin{longtable}{>{\raggedright}p{0.5cm}>{\raggedright}p{4.0cm}>{\raggedright}p{5.6cm}>{\raggedright}p{2.3cm}>{\raggedright\arraybackslash}p{1.7cm}}
\toprule
TT & Khẳng định & Thí nghiệm và điều đã đo & Nguồn & Trạng thái \\
\midrule\endfirsthead
\toprule
TT & Khẳng định & Thí nghiệm và điều đã đo & Nguồn & Trạng thái \\
\midrule\endhead
1 & Nơ-ron \nt{GABA} chiếm từ một phần năm đến một phần tư số nơ-ron vỏ não & Nhuộm miễn dịch GABA ở $10$ vùng vỏ não khỉ: khoảng $25\,\%$ nơ-ron ở $8$ vùng, dưới $20\,\%$ ở vỏ thị giác sơ cấp. Tổng quan về sự đa dạng của nơ-ron ức chế vỏ não & \cite{Hendry1987,Markram2004} & đã đo \\
2 & Ức chế làm gì phần lớn do vị trí đậu quyết định: đuôi gai, thân, chỗ sinh xung, đầu tận cùng của dây dẫn khác & Tổng quan: các lớp nơ-ron \nt{GABA} vỏ não khác nhau ở đích trên nơi nhận. Ghi đồng thời từ thân và đuôi gai của nơ-ron tháp: ức chế thuận mạnh hơn nhiều ở thân so với ở đuôi gai. Ức chế trên đầu tận cùng của dây dẫn khác ở tủy sống làm giảm lượng chất dẫn truyền được nhả của một đường vào (tổng quan các phép đo) & \cite{Markram2004,Pouille2001,Rudomin1999} & đã đo \\
3 & Đổi dấu qua trung gian tốn một độ trễ, khoảng một mili giây; ức chế theo sau kích thích thu hẹp cửa sổ trùng khớp & Nơ-ron tháp hồi hải mã chuột cống trong lát cắt: ức chế qua một trung gian đến theo sau kích thích trực tiếp với độ trễ ngắn, và cửa sổ trong đó hai đầu vào kích thích cộng lại tới ngưỡng nhỏ hơn $2\ms$. Ở đuôi gai, nơi ức chế này yếu hơn, cửa sổ rộng hơn & \cite{Pouille2001} & đã đo \\
4 & Phép cấm $a\wedge\bar b$~\eqref{eq:veto} cùng với OR và hằng một là đầy đủ về chức năng (mệnh đề~\ref{prop:gabacomplete}) & Chứng minh trong chương. Kết quả kinh điển: mạng phần tử ngưỡng có đầu vào ức chế hiện thực được mọi biểu thức logic. Khẳng định về mô hình, chưa có thí nghiệm & \cite{McCullochPitts1943} & lý thuyết \\
5 & Phép cấm có độ trễ cho bộ phát hiện hướng chuyển động với vận tốc tối ưu $v^*=\Delta x/d$~\eqref{eq:vstar} & Võng mạc thỏ: tính chọn lọc hướng được tạo ra bởi ức chế, dập tắt kích thích khi chuyển động theo hướng ``không''. Ghi riêng đầu vào kích thích và ức chế của các tế bào chọn lọc cho thấy tế bào amacrine hình sao (\nt{GABA}) ức chế chúng không đối xứng --- chỉ bằng các nhánh hướng về phía ``không''. Công thức $v^*$ là suy luận của chương & \cite{Barlow1965,Fried2002} & đã đo \\
6 & Ức chế sun chia điện áp chứ không trừ~\eqref{eq:shunt} (hình~\ref{fig:shunt}) & Định luật Ohm cho bộ chia gồm ba độ dẫn, với điện thế cân bằng của clo gần mức nghỉ; cho nơ-ron không phát xung & suy luận trong chương & lý thuyết \\
7 & Với tần số xung, sun tác động bằng phép trừ, còn phép chia có được nhờ ức chế cùng với nhiễu nền cân bằng & Mô hình: ở nơ-ron đang phát xung, dòng qua sun gần như không đổi, và hiệu ứng lên tần số là phép trừ. Thực nghiệm: đưa vào nơ-ron tháp vỏ cảm giác thân thể chuột cống (lát cắt) các độ dẫn nền kích thích và ức chế bằng kẹp động; tăng đồng thời và cân bằng cả hai làm chia độ dốc đáp ứng mà không dịch tần số đáng kể. Tổng quan: phép chia cho tổng hoạt động gặp ở nhiều vùng não & \cite{HoltKoch1997,Chance2002,Carandini2012} & đã đo \\
8 & Khi $\beta>1$, ức chế lẫn nhau chọn ra một kẻ thắng (mệnh đề~\ref{prop:wta}, \eqref{eq:wta}) & Các trị riêng $-(1\pm\beta)/\tau$ là suy luận trong chương. Tần số $0.73$ và $0.53$ khi $\beta=0.5$ là điểm bất động $r_{1,2}=(I_{1,2}-\beta I_{2,1})/(1-\beta^2)$; không có script trong \texttt{models/}. Chương không dẫn thí nghiệm trực tiếp & suy luận trong chương & lý thuyết \\
9 & Đặt lại bộ nhớ bằng tín hiệu qua \nt{GABA}; hai nhóm ức chế lẫn nhau là một mạch chốt & Suy ra từ hình~\ref{fig:bistable} và mệnh đề~\ref{prop:wta}. Chưa có thí nghiệm & suy luận trong chương & lý thuyết \\
10 & Trạng thái cân bằng của vòng \nt{GLUT}--\nt{GABA} ổn định nếu ức chế đủ mạnh và đủ nhanh (mệnh đề~\ref{prop:ei}) & Mô hình hai quần thể~\eqref{eq:ei}; các điều kiện (i) và (ii) là định thức và vết của ma trận, được suy ra trong chương & \cite{WilsonCowan1972} & lý thuyết \\
11 & Với $w_{EE}=1.5$ và $w_{EI}w_{IE}=2$, ức chế nhanh ($\tau_I=2\ms$) giữ $E^*=0.67h$, ức chế chậm ($30\ms$) gây dao động (hình~\ref{fig:ei}) & Nghiệm số của~\eqref{eq:ei}, script \texttt{figures/make\_ei.py} & tính toán & mô hình \\
12 & Vỏ não làm việc ở chế độ ổn định nhờ ức chế; dấu hiệu --- đẩy thêm nơ-ron \nt{GABA} thì tần số của chúng giảm~\eqref{eq:isn} & Lý thuyết đã dự đoán đáp ứng nghịch lý. Thực nghiệm: ghi nội bào ở vỏ thị giác sơ cấp của mèo --- khi đáp ứng bị kích thích xung quanh ức chế, độ dẫn ức chế không tăng mà giảm cùng với độ dẫn kích thích. Kích thích quang di truyền các nơ-ron PV của chuột nhắt ở vỏ thị giác, vỏ cảm giác thân thể và vỏ vận động làm giảm tần số của nơ-ron ức chế, kể cả khi không có kích thích và khi gây mê nhẹ. Hệ số $-1/3$ với số liệu hình~\ref{fig:ei} là suy luận của chương & \cite{Tsodyks1997isn,Ozeki2009,Sanzeni2020} & đã đo \\
13 & Vòng ``\nt{GLUT} kích thích \nt{GABA}, \nt{GABA} ức chế \nt{GLUT}'' sinh ra nhịp nhanh với chu kỳ $12$--$50\ms$ & Kích thích quang di truyền các nơ-ron \nt{GABA} nhanh ở vỏ cảm giác thân thể chuột nhắt với tần số $8$--$200$\,Hz tăng cường có chọn lọc nhịp gamma ($20$--$80$\,Hz, chu kỳ $12$--$50\ms$); kích thích nơ-ron \nt{GLUT} chỉ tăng cường các nhịp chậm hơn & \cite{Cardin2009,Buzsaki2012} & đã đo \\
\bottomrule
\end{longtable}}
